\documentclass[10pt,a4paper]{amsart}
\usepackage{amsmath,amssymb,amsthm,graphicx}
% Retained preamble for the Stacks Project chapter {CH} (shared project preamble).
% The chapter file pinned-original-source.tex only carries \input{{preamble}} on its first line,
% so no preamble text lives inside the chapter itself.  The Stacks Project uses one single shared
% preamble for every chapter: the snapshot stages it three times (sources/stacks-chow/preamble.tex,
% sources/stacks-etale-cohomology/preamble.tex and focused-seeds/AGP/sources/stacks-snapshot-20261011/
% preamble.tex) and those copies are identical in size.  The block below is copied mechanically,
% line for line, from sources/stacks-chow/preamble.tex lines 14-18 and 154-217: the Xy-pic setup and
% the native theorem declarations together with the mathematical macros.  Nothing was added, removed
% or redefined; the wrapper class and the amsmath/amssymb/amsthm/graphicx base packages are supplied
% by the assembler.
% For commutative diagrams we use Xy-pic
% Chapter: stacks-spaces-more-morphisms.tex -- More on Morphisms of Spaces.
\usepackage[all]{xy}

% We use 2cell for 2-commutative diagrams.
\xyoption{2cell}

% Theorem environments.
%
\theoremstyle{plain}
\newtheorem{theorem}[subsection]{Theorem}
\newtheorem{proposition}[subsection]{Proposition}
\newtheorem{lemma}[subsection]{Lemma}

\theoremstyle{definition}
\newtheorem{definition}[subsection]{Definition}
\newtheorem{example}[subsection]{Example}
\newtheorem{exercise}[subsection]{Exercise}
\newtheorem{situation}[subsection]{Situation}

\theoremstyle{remark}
\newtheorem{remark}[subsection]{Remark}
\newtheorem{remarks}[subsection]{Remarks}

\numberwithin{equation}{subsection}

% Macros
%
\def\lim{\mathop{\mathrm{lim}}\nolimits}
\def\colim{\mathop{\mathrm{colim}}\nolimits}
\def\Spec{\mathop{\mathrm{Spec}}}
\def\Hom{\mathop{\mathrm{Hom}}\nolimits}
\def\Ext{\mathop{\mathrm{Ext}}\nolimits}
\def\SheafHom{\mathop{\mathcal{H}\!\mathit{om}}\nolimits}
\def\SheafExt{\mathop{\mathcal{E}\!\mathit{xt}}\nolimits}
\def\Sch{\mathit{Sch}}
\def\Mor{\mathop{\mathrm{Mor}}\nolimits}
\def\Ob{\mathop{\mathrm{Ob}}\nolimits}
\def\Sh{\mathop{\mathit{Sh}}\nolimits}
\def\NL{\mathop{N\!L}\nolimits}
\def\CH{\mathop{\mathrm{CH}}\nolimits}
\def\proetale{{pro\text{-}\acute{e}tale}}
\def\etale{{\acute{e}tale}}
\def\QCoh{\mathit{QCoh}}
\def\Ker{\mathop{\mathrm{Ker}}}
\def\Im{\mathop{\mathrm{Im}}}
\def\Coker{\mathop{\mathrm{Coker}}}
\def\Coim{\mathop{\mathrm{Coim}}}

% Boxtimes
%
\DeclareMathSymbol{\boxtimes}{\mathbin}{AMSa}{"02}

%
% Macros for moduli stacks/spaces
%
\def\QCohstack{\mathcal{QC}\!\mathit{oh}}
\def\Cohstack{\mathcal{C}\!\mathit{oh}}
\def\Spacesstack{\mathcal{S}\!\mathit{paces}}
\def\Quotfunctor{\mathrm{Quot}}
\def\Hilbfunctor{\mathrm{Hilb}}
\def\Curvesstack{\mathcal{C}\!\mathit{urves}}
\def\Polarizedstack{\mathcal{P}\!\mathit{olarized}}
\def\Complexesstack{\mathcal{C}\!\mathit{omplexes}}
% \Pic is the operator that assigns to X its picard group, usage \Pic(X)
% \Picardstack_{X/B} denotes the Picard stack of X over B
% \Picardfunctor_{X/B} denotes the Picard functor of X over B
\def\Pic{\mathop{\mathrm{Pic}}\nolimits}
\def\Picardstack{\mathcal{P}\!\mathit{ic}}
\def\Picardfunctor{\mathrm{Pic}}
\def\Deformationcategory{\mathcal{D}\!\mathit{ef}}
\newenvironment{reference}{}{}
\newenvironment{slogan}{}{}
\newenvironment{history}{}{}
\usepackage{cleveref}
\crefname{originallemma}{original lemma}{original lemmas}
\Crefname{originallemma}{Original lemma}{Original lemmas}
\makeatletter
\newlabel{lemma-finite-after-blowing-up}{{0B4K}{source}{}{sourcepointer-1}{}}
\newlabel{lemma-finite-after-blowing-up@cref}{{[originallemma][0][]0B4K}{[1][0][]source}{}{}{}}
\makeatother
\crefname{originallemma}{original lemma}{original lemmas}
\Crefname{originallemma}{Original lemma}{Original lemmas}
\makeatletter
\newlabel{more-morphisms-lemma-ample-family-ample-relative}{{0GTP}{source}{}{sourcepointer-2}{}}
\newlabel{more-morphisms-lemma-ample-family-ample-relative@cref}{{[originallemma][0][]0GTP}{[1][0][]source}{}{}{}}
\makeatother
\crefname{originallemma}{original lemma}{original lemmas}
\Crefname{originallemma}{Original lemma}{Original lemmas}
\makeatletter
\newlabel{more-morphisms-lemma-blow-up-ample-family}{{0GTT}{source}{}{sourcepointer-3}{}}
\newlabel{more-morphisms-lemma-blow-up-ample-family@cref}{{[originallemma][0][]0GTT}{[1][0][]source}{}{}{}}
\makeatother
\crefname{originallemma}{original lemma}{original lemmas}
\Crefname{originallemma}{Original lemma}{Original lemmas}
\makeatletter
\newlabel{perfect-lemma-resolution-property-ample-family}{{0GMM}{source}{}{sourcepointer-4}{}}
\newlabel{perfect-lemma-resolution-property-ample-family@cref}{{[originallemma][0][]0GMM}{[1][0][]source}{}{}{}}
\makeatother
\crefname{originallemma}{original lemma}{original lemmas}
\Crefname{originallemma}{Original lemma}{Original lemmas}
\makeatletter
\newlabel{spaces-divisors-lemma-strict-transform}{{0864}{source}{}{sourcepointer-5}{}}
\newlabel{spaces-divisors-lemma-strict-transform@cref}{{[originallemma][0][]0864}{[1][0][]source}{}{}{}}
\makeatother
\crefname{originallemma}{original lemma}{original lemmas}
\Crefname{originallemma}{Original lemma}{Original lemmas}
\makeatletter
\newlabel{spaces-limits-lemma-there-is-a-scheme-finite-over-refined}{{0GUN}{source}{}{sourcepointer-6}{}}
\newlabel{spaces-limits-lemma-there-is-a-scheme-finite-over-refined@cref}{{[originallemma][0][]0GUN}{[1][0][]source}{}{}{}}
\makeatother
\crefname{originallemma}{original lemma}{original lemmas}
\Crefname{originallemma}{Original lemma}{Original lemmas}
\makeatletter
\newlabel{spaces-perfect-lemma-resolution-property-goes-down-finite-flat}{{0GUV}{source}{}{sourcepointer-7}{}}
\newlabel{spaces-perfect-lemma-resolution-property-goes-down-finite-flat@cref}{{[originallemma][0][]0GUV}{[1][0][]source}{}{}{}}
\makeatother
\title{Every qcqs space with finitely many irreducible components becomes resolution-property after a dense admissible blow-up}
\author{The Stacks Project Authors}
\date{stacks-project:spaces-more-morphisms.tex:local-snapshot-20261011; GNU FDL 1.2 or later}
\begin{document}
\maketitle
\noindent\textit{Editorial note.} Statement span is lines 13832-13839 of sources/stacks-spaces-more-morphisms/pinned-original-source.tex; proof span is lines 13841-13878, the complete author proof (36 body lines).\par\medskip
\noindent\textit{Editorial note.} Navigation stubs: 7 cited-but-not-included labels are supplied as non-load-bearing pointers with their pinned tags where the snapshot tag index has them.\par\medskip
\noindent\textit{Editorial note.} extra\_packages is empty on purpose: the retained preamble already loads xy with options and a second declaration clashes in the assembler wrapper.\par\medskip
\noindent\textit{Editorial note.} Text is retained verbatim, including the author's own typos and placeholders.\par\medskip
\noindent\textit{Editorial note.} Capstone of the spaces chapter's resolution-property series; the parenthetical (details omitted) covers only the harmless shrinking f\textasciicircum{}\{-1\}(U) subset V after shrinking U. Cross-chapter dependency on card L16087 is now satisfied (it assembled this morning).\par\medskip
\noindent\textit{Allowed standard prerequisites (editorial).}\begin{itemize}
\item Explicit prerequisite boundary. The chapter's finite-refinement limit theorem (Limits of Spaces), the schemes-side blow-up-ample-family theorem (card L16087, now on the shelf), the strict-transform finiteness lemma (this chapter's lemma finite-after-blowing-up, same wave L16118), projectivity of strict transforms (Divisors on Spaces), ample-family-with-relative-ample machinery (More on Morphisms), and the two resolution-property transfers (Derived Categories of Schemes/Spaces).
\end{itemize}\medskip
\ifdefined\hypertarget\hypertarget{sourcepointer-1}{}\fi
\noindent\textit{Source reference.} Original lemma 0B4K: sources/stacks-spaces-more-morphisms/pinned-original-source.tex, source lines 8913--8921. Cited by the retained statement or proof but not part of the collected spans; supplied only as a navigation pointer. The argument does not depend on its content.\par\medskip
\ifdefined\hypertarget\hypertarget{sourcepointer-2}{}\fi
\noindent\textit{Source reference.} Original lemma 0GTP: sources/stacks-spaces-more-morphisms/pinned-original-source.tex, source lines 13878--13878. Cited by the retained statement or proof but not part of the collected spans; supplied only as a navigation pointer. The argument does not depend on its content.\par\medskip
\ifdefined\hypertarget\hypertarget{sourcepointer-3}{}\fi
\noindent\textit{Source reference.} Original lemma 0GTT: sources/stacks-spaces-more-morphisms/pinned-original-source.tex, source lines 13878--13878. Cited by the retained statement or proof but not part of the collected spans; supplied only as a navigation pointer. The argument does not depend on its content.\par\medskip
\ifdefined\hypertarget\hypertarget{sourcepointer-4}{}\fi
\noindent\textit{Source reference.} Original lemma 0GMM: sources/stacks-spaces-more-morphisms/pinned-original-source.tex, source lines 13878--13878. Cited by the retained statement or proof but not part of the collected spans; supplied only as a navigation pointer. The argument does not depend on its content.\par\medskip
\ifdefined\hypertarget\hypertarget{sourcepointer-5}{}\fi
\noindent\textit{Source reference.} Original lemma 0864: sources/stacks-spaces-more-morphisms/pinned-original-source.tex, source lines 13878--13878. Cited by the retained statement or proof but not part of the collected spans; supplied only as a navigation pointer. The argument does not depend on its content.\par\medskip
\ifdefined\hypertarget\hypertarget{sourcepointer-6}{}\fi
\noindent\textit{Source reference.} Original lemma 0GUN: sources/stacks-spaces-more-morphisms/pinned-original-source.tex, source lines 13878--13878. Cited by the retained statement or proof but not part of the collected spans; supplied only as a navigation pointer. The argument does not depend on its content.\par\medskip
\ifdefined\hypertarget\hypertarget{sourcepointer-7}{}\fi
\noindent\textit{Source reference.} Original lemma 0GUV: sources/stacks-spaces-more-morphisms/pinned-original-source.tex, source lines 13878--13878. Cited by the retained statement or proof but not part of the collected spans; supplied only as a navigation pointer. The argument does not depend on its content.\par\medskip
\begin{lemma}
\label{lemma-blow-up-resolution-property}
Let $S$ be a scheme. Let $X$ be a quasi-compact and quasi-separated
algebraic space over $S$. Assume that $|X|$ has finitely many irreducible
components. There exists a dense quasi-compact open $U \subset X$
and a $U$-admissible blowing up $X' \to X$ such that the algebraic space
$X'$ has the resolution property.
\end{lemma}

\begin{proof}
By Limits of Spaces, Lemma
\ref{spaces-limits-lemma-there-is-a-scheme-finite-over-refined}
there exists a surjective, finite, and finitely presented morphism
$f : Y \to X$ where $Y$ is a scheme and a quasi-compact dense open
$U \subset X$ such that $f^{-1}(U) \to U$ is finite \'etale.
By More on Morphisms, Lemma \ref{more-morphisms-lemma-blow-up-ample-family}
there is a quasi-compact dense open $V \subset Y$ and a $V$-admissible
blowing up $Y' \to Y$ such that $Y'$ has an ample family of invertible
modules. After shrinking $U$ we may assume that $f^{-1}(U) \subset V$
(details omitted). Hence $f' : Y' \to X$ is finite \'etale over $U$
and in particular, the morphism $(f')^{-1}(U) \to U$ is finite locally free.
By Lemma \ref{lemma-finite-after-blowing-up} there is a $U$-admissible
blowing up $X' \to X$ such that the strict transform $Y''$ of $Y'$
is finite locally free over $X'$. Picture
$$
\xymatrix{
Y'' \ar[d] \ar[r]_g &
Y' \ar[r] &
Y \ar[d] \\
X' \ar[rr] & & X
}
$$
Since $g : Y'' \to Y'$ is a blowing up
(Divisors on Spaces, Lemma \ref{spaces-divisors-lemma-strict-transform})
in the inverse image of the center of $X' \to X$, we see that
$g : Y'' \to Y'$ is projective and that there exists some 
$g$-ample invertible module on $Y''$. Hence by
More on Morphisms, Lemma \ref{more-morphisms-lemma-ample-family-ample-relative}
we see that $Y''$ has an ample family of invertible modules.
Hence $Y''$ has the resolution property, see
Derived Categories of Schemes, Lemma
\ref{perfect-lemma-resolution-property-ample-family}.
We conclude that $X'$ has the resolution property
by
Derived Categories of Spaces, Lemma
\ref{spaces-perfect-lemma-resolution-property-goes-down-finite-flat}.
\end{proof}

\end{document}
